% ═══════════════════════════════════════════════════════════════
%  front/preface.tex
%  Предисловие
% ═══════════════════════════════════════════════════════════════

\chapter*{Предисловие}
\addcontentsline{toc}{chapter}{Предисловие}
\markboth{Предисловие}{Предисловие}

\section*{Зачем эта книга}

Я мечтал о ней двадцать лет.

Двадцать лет назад я впервые увидел, как мигает светодиод
на Arduino. Это было чудо: несколько строк кода~--- и вот
уже железо оживает, слушается, отвечает. Но чуда было мало.
Хотелось, чтобы Arduino <<разговаривала>> с компьютером, чтобы
на экране рисовался график, чтобы всё это работало \emph{под
Linux}, а не только в окошке Arduino IDE.

Тогда я не знал, как это сделать. Не знал, что такое
\texttt{termios}, \texttt{libudev}, \texttt{GIOChannel}.
Не знал, как соединить программу на C с графическим
интерфейсом на GTK. Информация была~--- но \emph{разбросана}
по форумам, мануалам и редким книгам на английском. Собрать
её в единую картину было \emph{очень} трудно.

Эта книга~--- попытка \emph{собрать} эту картину для вас.
Чтобы вы не тратили двадцать лет, а прошли путь за несколько
месяцев. Чтобы у вас было \emph{практическое руководство},
а не набор разрозненных статей.

\section*{Для кого эта книга}

Книга рассчитана на \textbf{программистов, знакомых с~C}, но
не знакомых с:

\begin{itemize}
    \item \textbf{GTK}~--- библиотекой для графических интерфейсов.
    \item \textbf{USB}~--- как он выглядит со стороны Linux.
    \item \textbf{termios}~--- настройкой последовательного порта.
    \item \textbf{udev}~--- идентификацией устройств в Linux.
\end{itemize}

Если вы никогда не работали с железом~--- не страшно.
Мы разберём всё с нуля. Если вы никогда не писали
GUI-приложений~--- тоже не страшно. GTK~--- это просто ещё
один API.

Что нужно знать \emph{обязательно}: C, Linux, терминал.

\section*{Как читать книгу}

Книга построена как \textbf{один проект}. Каждая глава~---
разбор одного модуля приложения \textbf{Arduino Control}.
В конце книги у вас будет \emph{работающее приложение} и
\emph{понимание}, как его написать с нуля.

Читайте \textbf{последовательно}, от главы к главе. Или~
--- \textbf{по оглавлению}, если ищете конкретную задачу.
Каждая глава \emph{самодостаточна}, но \emph{связана}
с остальными.

\begin{ideabox}
\textbf{Вся книга~--- это один проект.} Не набор
<<учебных примеров>>, а \emph{настоящая программа}~---
со всеми её граблями, странностями и инженерными
компромиссами.
\end{ideabox}

\section*{Структура книги}

Книга состоит из \textbf{пяти частей}:

\begin{enumerate}
    \item \textbf{Введение в проект}~--- обзор и архитектура.
    \item \textbf{Работа с железом}~--- порт, udev, выдёргивание USB.
    \item \textbf{Логика приложения}~--- буферы, конфиг, обработка.
    \item \textbf{Интерфейс}~--- GTK, виджеты, Cairo.
    \item \textbf{Завершение}~--- жизненный цикл, отладка.
\end{enumerate}

Плюс~--- \textbf{приложения}: ссылки на исходники, инструкция
по сборке, глоссарий.

\section*{Благодарности}

Прежде всего~--- \textbf{Arduino}. За то, что сделали
микроконтроллеры \emph{доступными}~--- и для школьника,
и для инженера. За открытую платформу, открытый код,
открытое сообщество.

\textbf{GNOME}~--- за GTK. За то, что можно писать
красивые интерфейсы \emph{на C}. За то, что GTK~---
не <<магия>>, а \emph{инструмент}, который можно
понять.

\textbf{Сообществу Linux}~--- за то, что вы \emph{есть}.
За Stack Overflow, за форумы, за man-страницы, за тысячи
часов ответов на вопросы, которые я задавал (и которые
задам ещё).

\textbf{Вам}, читателю~--- за то, что взяли эту книгу
в руки. Надеюсь, она будет полезна.

И~--- \textbf{AI-ассистенту Claude}. Без него эта книга
осталась бы \emph{замыслом}. Он помог \emph{структурировать}
материал, \emph{отредактировать} текст, \emph{подготовить}
примеры кода, \emph{найти} грабли и \emph{объяснить}
их просто. Это не <<замена автора>>, а \emph{партнёрство}:
человек + машина, где каждый делает то, что умеет лучше.

\begin{notebox}
\textbf{О коде в книге.} Код~--- \emph{не} единственно
возможный. Это \emph{наш} вариант, проверенный на практике.
Если вы видите \emph{другой} путь~--- попробуйте. Возможно,
он лучше.
\end{notebox}

\section*{Об исходниках}

Весь исходный код проекта \textbf{доступен} на GitHub.
Ссылка~--- в \textbf{приложении~А}. Сборка проекта~---
в \textbf{приложении~Б}. Глоссарий терминов~---
в \textbf{приложении~В}.

\vspace{1cm}

\begin{flushright}
\itshape
Александр Онацкий\\
\upshape\small 2027
\end{flushright}